\documentclass{beamer}

\mode<presentation>
\title{SoC1\&2: RISC-V Place \& Route with Innovus}
\author[AS]{Andreas Siggelkow \newline \texttt{andreas.siggelkow@rwu.de}}
\date[2026]{\today}
\subtitle{RWU}
\institute{Elektrotechnik und Informatik}

%%%%%%%%%%%%%%%%%%%%%%%%%%%%%%%%
% Logo
%%%%%%%%%%%%%%%%%%%%%%%%%%%%%%%%
\pgfdeclareimage[height=1cm]{hslogo}{./hs}
\logo{\pgfuseimage{hslogo}}
%%%%%%%%%%%%%%%%%%%%%%%%%%%%%%%%
% Macros
%%%%%%%%%%%%%%%%%%%%%%%%%%%%%%%%
\usepackage[T1]{fontenc}
\usepackage{lmodern}
\usepackage{textcomp}
\definecolor{links}{HTML}{2A1B81}
\hypersetup{colorlinks,linkcolor=,urlcolor=links}
\usepackage{graphicx}
\usepackage{ragged2e, tabularx}
\newcolumntype{Y}{>{\RaggedRight\arraybackslash}X}
\usetheme{Singapore}
\usecolortheme{seahorse}
\usefonttheme[onlylarge]{structuresmallcapsserif}
\usefonttheme[onlysmall]{structurebold}
\setbeamertemplate{footline}[frame number]
\usepackage{listings}
\newcommand{\astitle}{nix}
\newcommand{\assubtitle}{nix}
\newcommand{\asfolie}[3]{\begin{frame}\frametitle{#1}\begin{block}{#2}{#3}\end{block}\end{frame}}

% Commands and tool output are set with listings: dashes, $, _, \ and quotes
% appear exactly as they have to be typed.
\definecolor{lstbg}{HTML}{F2F2F7}
\definecolor{lstcomment}{HTML}{5A5A8A}
\lstset{
  basicstyle=\ttfamily\scriptsize,
  columns=fullflexible,
  keepspaces=true,
  upquote=true,
  breaklines=true,
  breakatwhitespace=false,
  backgroundcolor=\color{lstbg},
  frame=none,
  xleftmargin=2pt,
  xrightmargin=2pt,
  aboveskip=4pt,
  belowskip=4pt,
  showstringspaces=false
}
\lstdefinestyle{sh}{
  morecomment=[l]{\#},
  commentstyle=\color{lstcomment}\itshape
}
\lstdefinestyle{out}{basicstyle=\ttfamily\tiny}
\newcommand{\cmd}[1]{\texttt{#1}}

\begin{document}
%%%%%%%%%%%%%%%%%%%%%%%%%%%%%%%%
% First pages
%%%%%%%%%%%%%%%%%%%%%%%%%%%%%%%%
\setcounter{tocdepth}{1}

\begin{frame}
  \titlepage
\end{frame}

\begin{frame}
  \frametitle{Outline}
  \tableofcontents
\end{frame}

%%%%%%%%%%%%%%%%%%%%%%%%%%%%%%%%
% Overview
%%%%%%%%%%%%%%%%%%%%%%%%%%%%%%%%
\section{Overview}
\renewcommand{\astitle}{Overview}
\begin{frame}
\frametitle{\astitle}
\textcolor{red}{Overview}

\url{https://github.com/asiggel/rwu-rv64iV2.0}
\end{frame}

\renewcommand{\assubtitle}{What is Innovus?}
\begin{frame}
\frametitle{\astitle}
\begin{block}{\assubtitle}
  \footnotesize
  \begin{itemize}
    \item Cadence's place \& route tool: gate-level netlist $\rightarrow$ layout
    \item Third step of the Cadence flow for the X-FAB XH018 tape-out:
      Xcelium (RTL) $\rightarrow$ Genus (synthesis) $\rightarrow$ \textbf{Innovus (P\&R)}
    \item Genus decides \emph{which} cells, Innovus decides \emph{where} they
      go and \emph{how} they are wired
    \item This deck (step 03): four stages up to the routed layout
    \begin{itemize}
      \item \textbf{fp}: floorplan: die and core size, pad ring
      \item \textbf{place}: put every standard cell and macro onto the core area
      \item \textbf{cts}: clock tree synthesis: buffer tree from the clock pin
        to every flip-flop
      \item \textbf{route}: draw the wires on the metal layers
    \end{itemize}
    \item Controlled by Tcl, like Genus; reads the Genus database directly
    \item Design: RV\_NoPipeline, chip top \cmd{cc\_chip\_top}
  \end{itemize}
\end{block}
\end{frame}

\renewcommand{\assubtitle}{Inputs and Outputs}
\begin{frame}
\frametitle{\astitle}
\begin{block}{\assubtitle}
  \scriptsize
  \begin{tabularx}{\linewidth}{@{}lY@{}}
    \textbf{In}                    & \\\hline
    \cmd{build\_genus/innovus/cc\_chip\_top.db} & everything in one database:
                                     netlist, Liberty, LEF, MMMC, SDC\\[1pt]
    \cmd{02\_rv64\_floorplan\_pnr.tcl} & die size, pad ring, the P\&R steps\\[4pt]
    \textbf{Out}                   & \\\hline
    \cmd{cc\_chip\_top\_<stage>.db} & checkpoint after each stage
                                     (open it again in the GUI)\\[1pt]
    \cmd{*\_timing.rpt}            & timing after CTS and after routing\\[1pt]
    \cmd{*\_drc.rpt}               & design rule violations (spacing, width, \dots)\\[1pt]
    \cmd{*\_antenna.rpt}           & antenna violations\\[1pt]
    \cmd{*\_connectivity.rpt}      & open or shorted nets\\[1pt]
    \cmd{*\_summary.rpt}           & cell count, area, wire length per layer\\[1pt]
    \cmd{cc\_chip\_top\_routed.db} & input of the post-route step\\
  \end{tabularx}
  \vspace{4pt}
  \begin{itemize}
    \item Hand-off netlists, DEF and GDS come from the post-route step
      (\cmd{05\_DokuINNOVUSPost}), after the last cell change
  \end{itemize}
\end{block}
\end{frame}

\renewcommand{\assubtitle}{The Chip in Numbers}
\begin{frame}
\frametitle{\astitle}
\begin{block}{\assubtitle}
  \footnotesize
  \begin{tabularx}{\linewidth}{@{}lY@{}}
    die              & 2600\,$\times$\,2600\,\textmu m\\[2pt]
    core             & 2016\,$\times$\,2016\,\textmu m (rows), 35\,\textmu m
                       to the IO ring for the power ring\\[2pt]
    IO ring          & 257\,\textmu m deep (= height of a pad cell), all four sides\\[2pt]
    pads             & 26: inputs, outputs, bidirectional (GPIO), supply\\[2pt]
    corners          & 4 corner cells, gaps closed with IO filler cells\\[2pt]
    cell area (Genus)& about 3.04\,mm\textsuperscript{2} incl.\ SPRAM macros\\[2pt]
    metal layers     & 4 (XH018, MET1 \dots\ top metal)\\[2pt]
  \end{tabularx}
  \vspace{4pt}
  \begin{itemize}
    \item The chip is \emph{core limited}: the core decides the size,
      26 pads would fit on a much smaller ring (mostly filler)
    \item The I-Mem macro is the largest block: 16\,KiB instead of 32\,KiB
      would save about 0.8\,mm\textsuperscript{2} (open item TECH-03)
  \end{itemize}
\end{block}
\end{frame}

%%%%%%%%%%%%%%%%%%%%%%%%%%%%%%%%
% Setup
%%%%%%%%%%%%%%%%%%%%%%%%%%%%%%%%
\section{Setup}
\renewcommand{\astitle}{Setup}
\begin{frame}
\frametitle{\astitle}
\textcolor{red}{Setup}
\end{frame}

\renewcommand{\assubtitle}{Check the Environment}
\begin{frame}[fragile]
\frametitle{\astitle}
\begin{block}{\assubtitle}
  \begin{itemize}
    \item Work on the Cadence lab PC; the Cadence tools (release 2026/27)
      are already in the \cmd{PATH}
  \end{itemize}
% VERIFY-LAB: install path and version string of Innovus 2026/27
\begin{lstlisting}[style=sh]
which innovus
innovus -version
\end{lstlisting}
  \begin{itemize}
    \item Synthesis must have run before:
\begin{lstlisting}[style=sh]
./scripts/genus/genus_syn.sh -b       # -> build_genus/innovus/cc_chip_top.db
\end{lstlisting}
    \item \cmd{XFAB\_PATHS} must be set, as for Genus: the flow takes the
      site, pad, corner and IO filler cell names from it (NDA)
    \item GUI: on the lab PC or in a VNC session. Layout over \cmd{ssh -X}
      from home is very slow
  \end{itemize}
\end{block}
\end{frame}

\renewcommand{\assubtitle}{NDA and Licenses}
\begin{frame}[fragile]
\frametitle{\astitle}
\begin{block}{\assubtitle}
  \begin{itemize}
    \item The layout contains X-FAB cells and the X-FAB layer stack:
      \textbf{never commit} anything from \cmd{build\_innovus/},
      no screenshots of the layout outside the lab
      (\cmd{.gitignore} blocks the directory)
    \item Innovus needs a license as long as it runs, also an idle GUI:
      close it when you are done (\cmd{exit})
% VERIFY-LAB: feature name and number of Innovus licenses
\begin{lstlisting}[style=sh]
lmstat -c $CDS_LIC_FILE -a | grep -i innovus
\end{lstlisting}
    \item Run the flow once in batch mode and then \emph{look} at the
      checkpoints in the GUI (\cmd{-{}-open}): the flow does not have to
      run again for every look
  \end{itemize}
\end{block}
\end{frame}

\renewcommand{\assubtitle}{Scripts and Results}
\begin{frame}
\frametitle{\astitle}
\begin{block}{\assubtitle}
  \footnotesize
  In \cmd{RWU\_RV/RV\_NoPipeline}:\\[3pt]
  \begin{tabularx}{\linewidth}{@{}lY@{}}
    \cmd{scripts/innovus/innovus\_pnr.sh}   & P\&R with GUI (default) or batch,
                                              step by step, view checkpoints\\[2pt]
    \cmd{scripts/02\_rv64\_floorplan\_pnr.tcl} & the Innovus flow itself\\[2pt]
    \cmd{scripts/02a\_rv64\_power.tcl}    & macros + power grid (sourced)\\[2pt]
    \cmd{scripts/xh018\_mmmc.tcl}           & corners and views (came with the Genus DB)\\[2pt]
    \cmd{build\_genus/innovus/}             & input: hand-off from Genus\\[2pt]
    \cmd{build\_innovus/}                   & all results\\
  \end{tabularx}
  \vspace{4pt}
  \begin{itemize}
    \item \cmd{innovus\_pnr.sh} only \emph{calls} \cmd{02\_rv64\_floorplan\_pnr.tcl}
    \item $\Rightarrow$ floorplan or flow changes are made \emph{only} in
      the Tcl script
    \item This step ends with the routed layout. Post-route optimization,
      fillers, metal fill, hand-off and GDS: next Innovus step,
      \cmd{05\_DokuINNOVUSPost}
  \end{itemize}
\end{block}
\end{frame}

%%%%%%%%%%%%%%%%%%%%%%%%%%%%%%%%
% Batch
%%%%%%%%%%%%%%%%%%%%%%%%%%%%%%%%
\section{P\&R without GUI}
\renewcommand{\astitle}{P\&R without GUI}
\begin{frame}
\frametitle{\astitle}
\textcolor{red}{Place \& route in batch mode}
\end{frame}

\renewcommand{\assubtitle}{Run}
\begin{frame}[fragile]
\frametitle{\astitle}
\begin{block}{\assubtitle}
  \footnotesize
\begin{lstlisting}[style=sh,basicstyle=\ttfamily\small]
cd RWU_RV/RV_NoPipeline
./scripts/innovus/innovus_pnr.sh -b
\end{lstlisting}
% VERIFY-LAB: replace with the output of a real run
\begin{lstlisting}[style=out]
=== [innovus] 1/3 run 02_rv64_floorplan_pnr.tcl up to 'route' (batch) ...
=== [innovus] 2/3 checkpoints
  fp: cc_chip_top_fp.db
  place: cc_chip_top_placed.db
  cts: cc_chip_top_cts.db
  route: cc_chip_top_routed.db
=== [innovus] 3/3 result
  worst setup slack: ...
  drc:              ...
  antenna:          ...
  connectivity:     ...
  pg_connectivity:  ...
--- next: ./scripts/innovus/innovus_post.sh -b
PASS  place & route of cc_chip_top up to 'route'
\end{lstlisting}
  \begin{itemize}
    \item PASS = all stages ran (exit code 0). DRC and antenna violations
      are reported, not judged: read the numbers!
  \end{itemize}
\end{block}
\end{frame}

\renewcommand{\assubtitle}{Options}
\begin{frame}
\frametitle{\astitle}
\begin{block}{\assubtitle}
  \footnotesize
  \begin{tabularx}{\linewidth}{@{}lY@{}}
    \cmd{-b}, \cmd{-{}-batch}        & no GUI (default: Innovus GUI)\\[2pt]
    \cmd{-s}, \cmd{-{}-stop STAGE}   & stop after \cmd{fp}, \cmd{place},
                                       \cmd{cts} or \cmd{route}
                                       (default)\\[2pt]
    \cmd{-o}, \cmd{-{}-open STAGE}   & run nothing, open the checkpoint of
                                       \cmd{STAGE} in the GUI\\[2pt]
    \cmd{-c}, \cmd{-{}-clean}        & delete \cmd{build\_innovus/} first\\[2pt]
    \cmd{-h}, \cmd{-{}-help}         & help\\
  \end{tabularx}
  \vspace{6pt}
  \begin{itemize}
    \item Every run starts again from the Genus database; \cmd{-s} only
      decides where it ends
    \item A run deletes the old checkpoints and reports first, also those
      of the post-route step: what is in \cmd{build\_innovus/} always
      belongs to the last run
    \item \cmd{-b -s fp}: try a floorplan change in a few seconds
  \end{itemize}
\end{block}
\end{frame}

\renewcommand{\assubtitle}{What Is in \texttt{build\_innovus/}?}
\begin{frame}
\frametitle{\astitle}
\begin{block}{\assubtitle}
  \scriptsize
  \begin{tabularx}{\linewidth}{@{}lY@{}}
    \cmd{innovus} $\rightarrow$                & link to \cmd{build\_genus/innovus/} (input)\\[1pt]
    \cmd{run.tcl}                              & generated: \cmd{STOP\_AFTER}, \cmd{source} the flow,
                                                 checks, \cmd{exit} (batch)\\[1pt]
    \cmd{run.log}, \cmd{run.logv}, \cmd{run.cmd} & log, verbose log, commands\\[1pt]
    \cmd{cc\_chip\_top\_fp.db}                 & checkpoint: floorplan, pad ring,
                                                 macros, power grid\\[1pt]
    \cmd{cc\_chip\_top\_placed.db}             & checkpoint: placed\\[1pt]
    \cmd{cc\_chip\_top\_cts.db}                & checkpoint: clock tree\\[1pt]
    \cmd{cc\_chip\_top\_routed.db}             & checkpoint: routed\\[1pt]
    \cmd{*\_fp\_summary.rpt}, \cmd{*\_route\_summary.rpt} & design summary\\[1pt]
    \cmd{*\_placed\_area.rpt}, \cmd{*\_route\_area.rpt}   & area\\[1pt]
    \cmd{*\_cts\_timing.rpt}, \cmd{*\_route\_timing.rpt}  & worst paths\\[1pt]
    \cmd{*\_drc.rpt}, \cmd{*\_antenna.rpt}, \cmd{*\_connectivity.rpt} & physical checks\\[1pt]
    \cmd{*\_pg\_connectivity.rpt}             & power grid after \cmd{route\_special}\\
  \end{tabularx}
  \vspace{3pt}
  The post-route step (\cmd{05\_DokuINNOVUSPost}) adds its results to the
  same directory.
\end{block}
\end{frame}

\renewcommand{\assubtitle}{Reading the Results}
\begin{frame}
\frametitle{\astitle}
\begin{block}{\assubtitle}
  \footnotesize
  \begin{tabularx}{\linewidth}{@{}lY@{}}
    connectivity & \textbf{must} be clean: an open net is a broken chip\\[2pt]
    timing       & now with real wires and a real clock tree: more precise
                   than in Genus. Compare the slack with the Genus report\\[2pt]
    DRC          & spacing, width, enclosure \dots\ of the drawn metal.
                   For tape-out: 0\\[2pt]
    antenna      & long wires collect charge during manufacturing and
                   can destroy a gate oxide. Fix: diodes or a layer jump\\
  \end{tabularx}
  \vspace{4pt}
  \begin{itemize}
    \item State of the first run (November 2026, smoke test):
    \begin{itemize}
      \item connectivity clean, setup met, routing complete
      \item about 2630 DRC and 159 antenna violations
    \end{itemize}
    \item That run had no power grid and no core fillers (added since,
      not yet run); the pad order is still unchecked (TECH-06)
  \end{itemize}
\end{block}
\end{frame}

%%%%%%%%%%%%%%%%%%%%%%%%%%%%%%%%
% GUI
%%%%%%%%%%%%%%%%%%%%%%%%%%%%%%%%
\section{P\&R with GUI}
\renewcommand{\astitle}{P\&R with GUI}
\begin{frame}
\frametitle{\astitle}
\textcolor{red}{Watching the chip grow}
\end{frame}

\renewcommand{\assubtitle}{Three Ways to See the Layout}
\begin{frame}[fragile]
\frametitle{\astitle}
\begin{block}{\assubtitle}
  \footnotesize
\begin{lstlisting}[style=sh]
cd RWU_RV/RV_NoPipeline
./scripts/innovus/innovus_pnr.sh                 # 1. whole flow in the GUI
./scripts/innovus/innovus_pnr.sh -s place        # 2. stop after a stage
./scripts/innovus/innovus_pnr.sh -o cts          # 3. open a checkpoint
\end{lstlisting}
  \begin{enumerate}
    \item The window opens with the empty design; you watch the pad ring,
      the placement, the clock tree and the wires appear.
      The window stays open at the end
    \item The flow ends after \cmd{fp}, \cmd{place} or \cmd{cts}: you
      continue by hand in the console (next slides)
    \item Fastest: nothing runs, the GUI shows a stage of the last run.
      Best for comparing stages: \cmd{-o fp}, \cmd{-o place}, \dots
  \end{enumerate}
  \begin{itemize}
    \item The terminal is the Innovus console; \cmd{exit} closes Innovus
  \end{itemize}
\end{block}
\end{frame}

% VERIFY-LAB: view names, menus and keys of the Innovus 2026/27 GUI
\renewcommand{\assubtitle}{Finding Your Way in the Window}
\begin{frame}
\frametitle{\astitle}
\begin{block}{\assubtitle}
  \footnotesize
  \begin{itemize}
    \item \textbf{Three views} (buttons top right):
    \begin{itemize}
      \item \emph{Floorplan view}: die, core, pads, macros
      \item \emph{Amoeba view}: one coloured area per hierarchy
        (\cmd{cpu}, \cmd{cgu}, \cmd{as\_jtag}, \dots): where did a block end up?
      \item \emph{Physical view}: every cell and every wire
    \end{itemize}
    \item \textbf{Layer panel} (right): switch metal layers, cells, pins on
      and off. Only one metal layer: you see its preferred direction
    \item \textbf{Zoom}: mouse wheel or drag a box; key \cmd{f}: fit the whole die
    \item \textbf{Select}: click a cell or wire; \cmd{q} shows its attributes
      (name, master cell, net, layer)
    \item \textbf{Find}: \emph{Edit $\rightarrow$ Find/Select Object}, e.g.\ all
      instances of \cmd{core/cpu/*}
  \end{itemize}
\end{block}
\end{frame}

\renewcommand{\assubtitle}{What to Look at, Stage by Stage}
\begin{frame}
\frametitle{\astitle}
\begin{block}{\assubtitle}
  \footnotesize
  \begin{tabularx}{\linewidth}{@{}lY@{}}
    \cmd{fp}    & die outline, core box, pad ring with corners and fillers,
                  26 pads in netlist order. Rows for the standard cells\\[3pt]
    \cmd{place} & the two SPRAM macros, the standard cells in rows.
                  Amoeba view: is the CPU close to the memories?
                  Where are the pads they talk to?\\[3pt]
    \cmd{cts}   & clock tree buffers. Clock tree debugger:
                  \emph{Clock $\rightarrow$ CCOpt Clock Tree Debugger}: tree from
                  \cmd{clk\_i} through the CGU to every flip-flop\\[3pt]
    \cmd{route} & all wires; switch the layers one by one.
                  DRC markers: \emph{Tools $\rightarrow$ Violation Browser}\\
  \end{tabularx}
\end{block}
\end{frame}

\renewcommand{\assubtitle}{Step by Step in the Console}
\begin{frame}[fragile]
\frametitle{\astitle}
\begin{block}{\assubtitle}
  After \cmd{innovus\_pnr.sh -s fp}, in the Innovus console:
\begin{lstlisting}[style=sh]
place_design                         # stage place
report_area
create_ccopt_clock_tree_spec         # stage cts
ccopt_design
report_timing -max_paths 5
routeDesign                          # stage route
report_timing -max_paths 5
verify_connectivity -type all
verify_drc                           # markers appear in the layout
\end{lstlisting}
  \begin{itemize}
    \item The same commands, in the same order, as in
      \cmd{02\_rv64\_floorplan\_pnr.tcl}
    \item Watch the layout after each command
    \item Nothing is saved: \cmd{write\_db my.db} if you want to keep it
  \end{itemize}
\end{block}
\end{frame}

\renewcommand{\assubtitle}{Following a Timing Path}
\begin{frame}[fragile]
\frametitle{\astitle}
\begin{block}{\assubtitle}
  \footnotesize
\begin{lstlisting}[style=sh]
report_timing -max_paths 1           # worst path, as text
\end{lstlisting}
% VERIFY-LAB: menu name of the timing debugger
  \begin{itemize}
    \item \emph{Timing $\rightarrow$ Debug Timing}: list of paths;
      select one $\rightarrow$ it is highlighted in the layout
    \item You see where the start and end flip-flops are placed and how
      long the wire between them is
    \item Typical findings: a path crosses the whole die to a pad,
      or it runs into a macro pin on the far side
    \item Compare with the Genus report: Genus estimated the wires,
      Innovus has them
  \end{itemize}
\end{block}
\end{frame}

%%%%%%%%%%%%%%%%%%%%%%%%%%%%%%%%
% Under the hood
%%%%%%%%%%%%%%%%%%%%%%%%%%%%%%%%
\section{Under the Hood}
\renewcommand{\astitle}{Under the Hood}
\begin{frame}
\frametitle{\astitle}
\textcolor{red}{The flow by hand}
\end{frame}

\renewcommand{\assubtitle}{a) One Innovus Call}
\begin{frame}[fragile]
\frametitle{\astitle}
\begin{block}{\assubtitle}
  \footnotesize
\begin{lstlisting}[style=sh]
cd RWU_RV/RV_NoPipeline
mkdir -p build_innovus && cd build_innovus
ln -sfn ../build_genus/innovus innovus
printf 'source ../scripts/02_rv64_floorplan_pnr.tcl\nexit\n' > run.tcl
innovus -batch -files run.tcl -log run -overwrite < /dev/null
\end{lstlisting}
  \begin{itemize}
    \item The flow reads \cmd{innovus/cc\_chip\_top.db} relative to the run
      directory: hence the link to the Genus results
    \item \cmd{-batch}: no window; without it the GUI opens
    \item Stop point: \cmd{set STOP\_AFTER place} before the \cmd{source};
      the flow then ends after writing \cmd{cc\_chip\_top\_placed.db}
    \item \cmd{-log run}: \cmd{run.log}, \cmd{run.logv}, \cmd{run.cmd};
      \cmd{-overwrite}: no numbered copies
  \end{itemize}
\end{block}
\end{frame}

\renewcommand{\assubtitle}{b) Floorplan}
\begin{frame}[fragile]
\frametitle{\astitle}
\begin{block}{\assubtitle}
  \footnotesize
\begin{lstlisting}[style=sh]
read_db innovus/${DESIGN_TOP}.db
set CORE_SIZE [expr {2600 - 2*257 - 2*35}]      ;# 2016 um
floorPlan -site $CORE_SITE -s $CORE_SIZE $CORE_SIZE 35 35 35 35
\end{lstlisting}
  \begin{itemize}
    \item \cmd{-site}: the row type of the high-density standard cells
      (name from \cmd{\$XFAB\_PATHS}, like all X-FAB cell names)
    \item \cmd{-s}: size of the \emph{core}. On this site Innovus keeps a
      margin of 257\,\textmu m on every side by itself: exactly the IO ring
    \item Result: die 2600\,$\times$\,2600\,\textmu m with standard-cell rows
      in the core
    \item Found by trying: margin 0 gives a core box from 257 to 2343\,\textmu m.
      The 35\,\textmu m margin on top holds the core power ring
      (VERIFY-LAB: the script prints die and core box)
  \end{itemize}
\end{block}
\end{frame}

\renewcommand{\assubtitle}{c) Pad Ring}
\begin{frame}
\frametitle{\astitle}
\begin{block}{\assubtitle}
  \footnotesize
  \begin{itemize}
    \item The script collects all pad instances of the netlist
      (by their pad cell type) and expects exactly 26
    \item Distribution: blocks of 7/7/7/5 pads, bottom $\rightarrow$ right
      $\rightarrow$ top $\rightarrow$ left, 60\,\textmu m pitch,
      fixed (\cmd{placeInstance \dots\ -fixed}), rotated to face outwards
    \item Corner cells: not in the netlist, added as physical-only
      instances (\cmd{addInst -physical})
    \item IO fillers (\cmd{addIoFiller}): close the gaps between the pads,
      so that the supply rails in the ring are continuous
    \item \textbf{First pass}: the order is the netlist order. The real
      assignment (which pad where, supply pads, ESD, simultaneous switching)
      comes from the IO databook (TECH-06)
  \end{itemize}
\end{block}
\end{frame}

\renewcommand{\assubtitle}{d) Placement, CTS, Routing}
\begin{frame}[fragile]
\frametitle{\astitle}
\begin{block}{\assubtitle}
  \footnotesize
\begin{lstlisting}[style=sh]
place_design                     ;  write_db cc_chip_top_placed.db
create_ccopt_clock_tree_spec
ccopt_design                     ;  write_db cc_chip_top_cts.db
routeDesign                      ;  write_db cc_chip_top_routed.db
\end{lstlisting}
  \begin{itemize}
    \item \cmd{place\_design}: places the standard cells (the macros are
      already fixed), optimizes timing on the way
    \item \cmd{create\_ccopt\_clock\_tree\_spec}: derives the clock trees from
      the SDC (\cmd{clk\_i}, generated \cmd{clk\_core}, \cmd{tck\_i})
    \item \cmd{ccopt\_design}: builds and balances the trees; after that the
      clocks are \emph{propagated}: real delays instead of ideal clocks
    \item \cmd{routeDesign}: global and detailed routing of all signal nets
    \item Before \cmd{routeDesign}: antenna fixing switched on (see g)
    \item Then the next Innovus step: post-route optimization
      (\cmd{05\_DokuINNOVUSPost})
  \end{itemize}
\end{block}
\end{frame}

\renewcommand{\assubtitle}{e) SRAM Macros (\texttt{02a\_rv64\_power.tcl})}
\begin{frame}
\frametitle{\astitle}
\begin{block}{\assubtitle}
  \footnotesize
  \begin{itemize}
    \item Placed \emph{before} the power grid and fixed, so that stripes
      and row rails are built around them
    \item I-Mem in the lower left corner of the core, D-Mem in the lower
      right; if both do not fit side by side, the D-Mem goes to the upper
      right. Sizes come from the macro LEF (\cmd{.base\_cell.bbox})
    \item Halo of 20\,\textmu m around each macro: no standard cells
      there, room for the supply connections
    \item The log line \cmd{AS-pnr: I-Mem WxH at (x,y)} shows the result
  \end{itemize}
\end{block}
\end{frame}

\renewcommand{\assubtitle}{f) Power Grid: Structure}
\begin{frame}
\frametitle{\astitle}
\begin{block}{\assubtitle}
  \footnotesize
  \begin{tabularx}{\linewidth}{@{}lY@{}}
    \textbf{part} & \textbf{what}\\\hline
    global connect & every supply pin of every cell to its net:
                     \cmd{VDD}/\cmd{GND} (core 1.8\,V), \cmd{VDDO}/\cmd{GNDO}
                     (IO 3.3\,V)\\
    core ring      & \cmd{VDD}/\cmd{GND} around the core, 12\,\textmu m wide;
                     horizontal sides \cmd{PG\_LAYER\_H}, vertical
                     \cmd{PG\_LAYER\_V} (top metal)\\
    stripes        & vertical \cmd{VDD}/\cmd{GND} pairs on the top metal,
                     every 200\,\textmu m, 8\,\textmu m wide\\
    rails          & \cmd{route\_special}: the row rails (followpins) on
                     metal 1, the supply pads to the ring, the macro supply
                     pins to ring and stripes\\
    IO supplies    & \cmd{VDDO}/\cmd{GNDO} only in the pad ring: by abutment of
                     pads, corners and IO fillers\\
  \end{tabularx}
  \vspace{4pt}
  \begin{itemize}
    \item Widths and pitch are a first estimate (a few mA at 1.25\,MHz);
      an IR-drop analysis (Voltus) is still open
  \end{itemize}
\end{block}
\end{frame}

\renewcommand{\assubtitle}{f) Power Grid: Commands}
\begin{frame}[fragile]
\frametitle{\astitle}
\begin{block}{\assubtitle}
\begin{lstlisting}[style=sh]
foreach p $PG_PINS_VDD {
  connect_global_net VDD -type pg_pin -pin_base_name $p -inst_base_name *
}                                          ;# same for GND, VDDO, GNDO
connect_global_net VDD -type tie_hi -inst_base_name *
add_rings -nets {VDD GND} -type core_rings -follow core \
  -layer [list top $PG_LAYER_H bottom $PG_LAYER_H \
               left $PG_LAYER_V right $PG_LAYER_V] \
  -width 12 -spacing 3 -offset 3
add_stripes -nets {VDD GND} -layer $PG_LAYER_V -direction vertical \
  -width 8 -spacing 3 -set_to_set_distance 200 -start_offset 100
route_special -nets {VDD GND} -connect {core_pin pad_pin block_pin}
\end{lstlisting}
  \scriptsize
  \begin{itemize}
    \item Optional before the ring: \cmd{add\_well\_taps},
      \cmd{add\_endcaps}, if \cmd{WELLTAP\_CELL}/\cmd{ENDCAP\_CELL} are set
    \item Check: \cmd{cc\_chip\_top\_pg\_connectivity.rpt}, and in the GUI
      the layers of the ring and the stripes
  \end{itemize}
\end{block}
\end{frame}

\renewcommand{\assubtitle}{g) Antenna Fixing}
\begin{frame}[fragile]
\frametitle{\astitle}
\begin{block}{\assubtitle}
\begin{lstlisting}[style=sh]
set_db route_design_detail_fix_antenna      true
set_db route_design_antenna_diode_insertion true
set_db route_design_antenna_cell_name       $ANTENNA_CELL
routeDesign
\end{lstlisting}
  \begin{itemize}
    \item During manufacturing a long wire collects charge before the
      upper metal connects it to a diffusion; the charge can break the
      gate oxide it ends on
    \item Rule: wire area / gate area below a limit per layer (LEF
      \cmd{ANTENNA*} data, X-FAB antenna deck)
    \item Fix 1: change to a higher layer near the gate (short pieces)
    \item Fix 2: an antenna diode next to the gate drains the charge
    \item Check: \cmd{cc\_chip\_top\_antenna.rpt}; sign-off is the antenna
      deck in \cmd{scripts/pv}
  \end{itemize}
\end{block}
\end{frame}

\renewcommand{\assubtitle}{h) Physical Checks}
\begin{frame}[fragile]
\frametitle{\astitle}
\begin{block}{\assubtitle}
  \footnotesize
  Added by \cmd{innovus\_pnr.sh} after routing:
\begin{lstlisting}[style=sh]
verify_connectivity -type all > cc_chip_top_connectivity.rpt   # in the flow
verify_drc -limit 100000 -report cc_chip_top_drc.rpt
verifyProcessAntenna -report cc_chip_top_antenna.rpt
\end{lstlisting}
  \begin{itemize}
    \item \cmd{-limit}: by default the DRC stops after a few hundred
      violations; the first pass has more
    \item These are Innovus' own checks with the LEF rules. Signoff for
      tape-out is done later with the X-FAB rule deck (\cmd{08\_DokuPV})
    \item The post-route step runs them again on the final layout
  \end{itemize}
\end{block}
\end{frame}

%%%%%%%%%%%%%%%%%%%%%%%%%%%%%%%%
% Workflow
%%%%%%%%%%%%%%%%%%%%%%%%%%%%%%%%
\section{Workflow}
\renewcommand{\astitle}{Workflow}
\begin{frame}
\frametitle{\astitle}
\textcolor{red}{Workflow}
\end{frame}

\renewcommand{\assubtitle}{From RTL to Layout}
\begin{frame}[fragile]
\frametitle{\astitle}
\begin{block}{\assubtitle}
\begin{lstlisting}[style=sh]
cd RWU_RV/RV_NoPipeline
./scripts/xcelium/xrun_regression.sh -c     # 01 does it work?
./scripts/genus/genus_syn.sh -b -c          # 02 gates, area, timing
./scripts/innovus/innovus_pnr.sh -b -c      # 03 layout up to routing
./scripts/innovus/innovus_pnr.sh -o route   #    look at it
./scripts/gatesim/gatesim.sh ...            # 04 gate-level sim
./scripts/innovus/innovus_post.sh -b        # 05 post-route, GDS
\end{lstlisting}
  \begin{itemize}
    \item Numbers = HowTo folders. Each step reads the result of the step
      before; after an RTL change all of them run again
    \item Batch mode, the GUI only for looking: saves
      licenses and time
  \end{itemize}
\end{block}
\end{frame}

\renewcommand{\assubtitle}{Changing the Floorplan}
\begin{frame}[fragile]
\frametitle{\astitle}
\begin{block}{\assubtitle}
  \footnotesize
  \begin{enumerate}
    \item Change \cmd{02\_rv64\_floorplan\_pnr.tcl}, e.g.\ \cmd{CORE\_SIZE}
    \item Quick check, floorplan only:
\begin{lstlisting}[style=sh]
./scripts/innovus/innovus_pnr.sh -b -s fp
./scripts/innovus/innovus_pnr.sh -o fp
\end{lstlisting}
    \item Then the whole flow: \cmd{innovus\_pnr.sh -b}
    \item Compare: area utilization (\cmd{\_placed\_area.rpt}), slack,
      DRC count
  \end{enumerate}
  \vspace{4pt}
  \begin{itemize}
    \item A smaller core is cheaper but harder to route: DRC violations
      grow when the core gets too full
    \item Die size and pad count are fixed in the spec (section 4.1.3):
      change the spec first, then the script
  \end{itemize}
\end{block}
\end{frame}

%%%%%%%%%%%%%%%%%%%%%%%%%%%%%%%%
% Pitfalls
%%%%%%%%%%%%%%%%%%%%%%%%%%%%%%%%
\section{Pitfalls}
\renewcommand{\astitle}{Pitfalls}
\begin{frame}
\frametitle{\astitle}
\textcolor{red}{Pitfalls}
\end{frame}

\renewcommand{\assubtitle}{No Genus Database}
\begin{frame}[fragile]
\frametitle{\astitle}
\begin{block}{\assubtitle}
\begin{lstlisting}
ERROR: .../build_genus/innovus/cc_chip_top.db not found
       run ./scripts/genus/genus_syn.sh -b first
\end{lstlisting}
  \begin{itemize}
    \item Synthesis has not run, or \cmd{build\_genus/} was deleted
    \item Genus results elsewhere: \cmd{GENUS\_ROOT=/path/to/genus/run ./scripts/innovus/innovus\_pnr.sh}
    \item X-FAB file not found while reading the database: the
      \cmd{\$XFAB\_PATHS} files have moved since the synthesis
      $\rightarrow$ run the synthesis again
  \end{itemize}
\end{block}
\end{frame}

\renewcommand{\assubtitle}{Pad Count}
\begin{frame}[fragile]
\frametitle{\astitle}
\begin{block}{\assubtitle}
\begin{lstlisting}
AS-pnr: expected 26 pad instances, found 24
\end{lstlisting}
  \begin{itemize}
    \item A port was added to or removed from \cmd{cc\_chip\_top},
      or Genus ran without \cmd{XFAB\_PADS}
    \item The pad ring is planned for exactly 26 pads; a new pad changes
      the ring: update the spec, the expected number and the pad list in
      \cmd{02\_rv64\_floorplan\_pnr.tcl}
  \end{itemize}
\end{block}
\end{frame}

\renewcommand{\assubtitle}{GUI, Display and Licenses}
\begin{frame}
\frametitle{\astitle}
\begin{block}{\assubtitle}
  \begin{itemize}
    \item The GUI needs an X display; \cmd{innovus\_pnr.sh} stops with an
      error if \cmd{\$DISPLAY} is empty
    \item A layout with thousands of cells over \cmd{ssh -X} is unusable:
      VNC or the lab PC
    \item The GUI keeps its license until \cmd{exit}; a forgotten window
      blocks a colleague
    \item Do not run two flows in the same \cmd{build\_innovus/} at the same time
    \item \cmd{-o} shows the \emph{last} run: after a \cmd{-s fp} run, the
      later checkpoints are gone
  \end{itemize}
\end{block}
\end{frame}

%%%%%%%%%%%%%%%%%%%%%%%%%%%%%%%%
% Q&A
%%%%%%%%%%%%%%%%%%%%%%%%%%%%%%%%
\section{Q\&A}
\renewcommand{\astitle}{Q\&A}
\begin{frame}
\frametitle{\astitle}
\vfill
\centering
\begin{Huge}\textcolor{red}{Questions?}\end{Huge}
\vfill
\end{frame}

\end{document}
